\section{Subharmonic limits of tilted norms}\label{sec:profiles}

We record exponential growth through logarithms of squared local
$L^2$ norms.  Pointwise logarithms are singular at zeros; the compactness
statement below instead controls masses on all compact and open sets,
which is what later comparisons of norms require.  Its compatibility
with affine weights will allow us to change the tilt as we subdivide.
Throughout the planar arguments,
norms and integrals are Euclidean and are not normalized by area.
We write $\D_r(a)$ for the open Euclidean disk of radius $r$ centered
at $a$, and $\D_r=\D_r(0)$; a bar denotes closure. For a measurable
planar set $E$, the notation $|E|$ denotes its Euclidean area.

\begin{lemma}[Extraction of a logarithmic profile]\label{lem:profile-extraction}
Let $X$ be an open subset of $\R^2$ or a compact smooth surface.  Suppose
that $S_j\to\infty$ and that $\mu_j$ are nonnegative Borel measures on
$X$ with $\mu_j(X)\le 1$.  There are a subsequence and an upper
semicontinuous function $V\colon X\to[-\infty,0]$ such that, for every
compact $F\subset X$ and every open $G\subset X$,
\begin{equation}\label{eq:profile-bounds}
 \limsup_{j\to\infty}\frac{\log\mu_j(F)}{2S_j}\le\sup_F V,
 \qquad
 \liminf_{j\to\infty}\frac{\log\mu_j(G)}{2S_j}\ge\sup_G V.
\end{equation}
Here $\log 0=\sup\varnothing=-\infty$.  More generally, if continuous
real functions $f_j$ converge locally uniformly to $f$, then
\[
 \begin{aligned}
 \limsup_{j\to\infty}\frac{1}{2S_j}
       \log\int_F e^{2S_j f_j}\,d\mu_j&\le\sup_F(V+f),\\
 \liminf_{j\to\infty}\frac{1}{2S_j}
       \log\int_G e^{2S_j f_j}\,d\mu_j&\ge\sup_G(V+f).
 \end{aligned}
\]
The same subsequence works for all these sets and weights.
\end{lemma}

\begin{proof}
Choose a countable basis $\{G_i\}_{i\ge1}$ consisting of relatively
compact open sets.  A diagonal subsequence makes all the limits
\[
 \ell_i=\lim_{j\to\infty}\frac{\log\mu_j(G_i)}{2S_j}
        \in[-\infty,0]
\]
exist.  Define $V(x)=\inf_{i:x\in G_i}\ell_i$.  If $V(x)<a$, one of
these basis neighborhoods has $\ell_i<a$, and $V<a$ throughout that
neighborhood.  Thus $V$ is upper semicontinuous.

For any finite $a>\sup_F V$, every point of $F$ has a basis
neighborhood with exponent less than $a$.  A finite subcover and
$\mu_j(F)\le\sum_i\mu_j(G_i)$ give the compact upper bound by $a$.
Letting $a\downarrow\sup_F V$ proves the first inequality, also when
this supremum is $-\infty$.  If $x\in G$, a basis neighborhood
$G_i$ with $x\in G_i\subset G$ gives a lower bound by
$\ell_i\ge V(x)$, proving the second inequality.

For the weighted upper bound, take $a>\sup_F(V+f)$ and choose the
neighborhood at each point small enough that its limiting measure
exponent plus the supremum of $f$ on it is less than $a$.  This is
possible by continuity of $f$ and the definition of $V$, including
at points where $V=-\infty$.  A finite cover and uniform convergence
of $f_j$ on its compact closure give the claimed upper bound.  For
the lower bound at a point $x\in G$ with $V(x)>-\infty$, choose
$G_i$ compactly contained in $G$ so small that
$f\ge f(x)-\eps$ there.  The weighted integral over $G_i$ has lower
exponent at least $\ell_i+f(x)-\eps$.  Let $\eps\downarrow0$ and
then take the supremum over $x$.
\end{proof}

We call $V$ a \emph{logarithmic profile}.  Multiplying the measures
by positive densities uniformly bounded above and below on compact
sets does not change the profile bounds: the logarithms of these
bounds disappear after division by $2S_j$.  We may also use the
compact upper bound for a relatively compact set by taking its
closure.  Normalizing mass on an open disk alone does not assert
that the profile has a finite value inside the disk; the
identically $-\infty$ case remains possible.

\begin{theorem}[Subharmonicity of logarithmic profiles]
\label{thm:subharmonic-profile}
On a fixed planar disk $\Omega$, let real smooth functions $U_j,p_j$
satisfy
\[
 \Delta U_j+K_j^2p_jU_j=0,
\]
where $K_j\ge0$.  Let $S_j>0$, $B_j\in\R^2$, and
$D_j=S_j+|B_j|$.  Assume
\begin{equation}\label{eq:profile-hypotheses}
 \begin{gathered}
 S_j\longrightarrow\infty,\qquad
 \frac{K_j}{D_j}\longrightarrow0,\\
 \sup_j\|p_j\|_\infty<\infty,\qquad
 \frac{K_j^2\|\nabla p_j\|_\infty}{S_jD_j}\longrightarrow0.
 \end{gathered}
\end{equation}
Local versions of the coefficient assumptions on relatively compact
subdisks suffice.  For positive constants $c_j$, suppose the measures
\[
 d\mu_j=c_j e^{-2B_j\cdot x}|U_j(x)|^2\,dx
\]
have mass at most one on $\Omega$.  Every upper semicontinuous
logarithmic profile with the compact, open, and weighted bounds of
Lemma~\ref{lem:profile-extraction} is subharmonic on $\Omega$,
with the identically $-\infty$ function allowed.
\end{theorem}

The main estimate concerns weights whose normalized gradient is
nonzero.  It is essential that the next lemma allows $D_j/S_j$ to
be arbitrarily large.

\begin{lemma}[A local Carleman estimate]\label{lem:carleman}
Assume the coefficient and parameter hypotheses of
Theorem~\ref{thm:subharmonic-profile}.  After passing to a subsequence,
write
\[
 S_j/D_j\to d,\qquad B_j/D_j\to b,
 \qquad d\ge0,\quad d+|b|=1.
\]
Let $\chi$ be a real smooth function near $x_0\in\Omega$ such that
$\Delta\chi(x_0)<0$ and $b+d\nabla\chi(x_0)\ne0$.  Put
$T_j(x)=-B_j\cdot x-S_j\chi(x)$.  There are a neighborhood
$\Omega_0$ of $x_0$, a constant $c>0$, and an index $j_0$ such that
\begin{equation}\label{eq:carleman}
 \|e^{T_j}(\Delta+K_j^2p_j)w\|_2
 \ge c\sqrt{S_j}\,D_j\,\|e^{T_j}w\|_2
\end{equation}
for all $j\ge j_0$ and all real $w\in C_c^\infty(\Omega_0)$.
\end{lemma}

The two parameters reflect different derivatives of the weight:
$\nabla T_j=-B_j-S_j\nabla\chi$ has scale $D_j$, whereas
$T_j''=-S_j\chi''$ has scale $S_j$.  The affine tilt can therefore
enlarge the gradient without enlarging the curvature.

\begin{proof}
Suppress the index and write $v=e^T w$.  The conjugated operator is
$P_++P_-$, where
\[
 P_+=\Delta+|\nabla T|^2+K^2p,
 \qquad P_-=-2\nabla T\cdot\nabla-\Delta T.
\]
On compactly supported functions, $P_+$ is symmetric and $P_-$ is
skew-symmetric.  Integration by parts gives
\begin{equation}\label{eq:commutator}
 \begin{split}
 2\langle P_+v,P_-v\rangle
  ={}&4\int T''(\nabla v,\nabla v)\\
    &+\int\bigl(4T''(\nabla T,\nabla T)-\Delta^2T
             +2K^2\nabla p\cdot\nabla T\bigr)v^2.
 \end{split}
\end{equation}
For clarity, the contribution from $\Delta$ is
$4\int T''(\nabla v,\nabla v)-\int\Delta^2T\,v^2$.
Multiplication by a real function $a$ contributes
$2\int\nabla a\cdot\nabla T\,v^2$; applying this with
$a=|\nabla T|^2+K^2p$ gives the remaining terms in
Equation~\eqref{eq:commutator}.

If Equation~\eqref{eq:carleman} failed on every sufficiently small
neighborhood, a diagonal choice would give a subsequence and real
smooth $v_j$ with supports shrinking to $x_0$, such that
\[
 \|v_j\|_2=1,
 \qquad
 \|(P_++P_-)v_j\|_2=o(\sqrt{S_j}D_j).
\]
All constants below are uniform on one fixed small neighborhood of
$x_0$.  There, $|\nabla T|\le CD$, $|T''|+|\Delta^2T|\le CS$,
and the last hypothesis in Equation~\eqref{eq:profile-hypotheses}
controls the potential term in Equation~\eqref{eq:commutator}.
Set $X=\|P_+v\|_2$, $Y=\|P_-v\|_2$, and
$G=\|\nabla v\|_2$.  Expanding the squared norm and using the
identity gives
\[
 X^2+Y^2\le CS(G^2+D^2).
\]
Testing $P_+$ against $v$ gives $G^2\le CD^2+X$, since $K=o(D)$.
Combining the two inequalities first yields
$X\le C(S+\sqrt S D)$ and then
\[
 G\le CD,
 \qquad X+Y\le C\sqrt S D.
\]
Here $D\ge S\to\infty$ absorbs the term $S$.

Uniformly on the shrinking supports,
\[
 \nabla T_j/D_j\to a=-b-d\nabla\chi(x_0)\ne0,
 \qquad T_j''/S_j\to H=-\chi''(x_0).
\]
The preceding bounds give
\[
 \frac{\sqrt{S_j}}{D_j}\le\frac1{\sqrt{S_j}}\longrightarrow0.
\]
Dividing the identity obtained by testing $P_+$ against $v_j$ by
$D_j^2$ therefore shows
\[
 \int|\nabla v_j/D_j|^2\longrightarrow |a|^2.
\]
Likewise,
\[
 \left\|\frac{\nabla T_j}{D_j}\cdot
                \frac{\nabla v_j}{D_j}\right\|_2
 \le \frac{\|P_-v_j\|_2+\|\Delta T_j\,v_j\|_2}{2D_j^2}
 \longrightarrow0.
\]
Replacing $\nabla T_j/D_j$ by its uniform limit consequently gives
$\|a\cdot\nabla v_j/D_j\|_2\to0$.

Let $e$ be a unit vector perpendicular to $a$.  In dimension two,
the energy perpendicular to $a$ has only this one direction.
Decomposing $\nabla v_j/D_j$ in the orthonormal basis
$\{a/|a|,e\}$, the parallel component tends to zero in $L^2$,
whereas the squared norm of the perpendicular component tends to
$|a|^2$.  Hence the first term in
Equation~\eqref{eq:commutator}, divided by $S_jD_j^2$, tends to
$4|a|^2H(e,e)$.  The other main Hessian term tends to $4H(a,a)$.
The two remaining terms tend to zero: their absolute values after
this division are bounded by
\[
 \frac{C}{D_j^2}
 \quad\hbox{and}\quad
 \frac{CK_j^2\|\nabla p_j\|_\infty}{S_jD_j},
\]
respectively.  We have thus proved
\[
 \begin{aligned}
 \frac{2\langle P_+v_j,P_-v_j\rangle}{S_jD_j^2}
 &\longrightarrow 4\bigl(|a|^2H(e,e)+H(a,a)\bigr)\\
 &=4|a|^2\operatorname{tr}H>0.
 \end{aligned}
\]
The squared norms of $P_+v_j$ and $P_-v_j$ are nonnegative, so this
contradicts the assumed smallness of $\|(P_++P_-)v_j\|_2^2$.
\end{proof}

We also need to remove one possible exceptional gradient from the
upper-test criterion.  The following elementary perturbation does
so without any initial integrability assumption.  An upper test at
a finite point $x_0$ is a smooth function agreeing with the tested
function at $x_0$ and dominating it in a neighborhood of $x_0$.
This is an upper-semicontinuous subharmonic variant of the standard
viscosity exceptional-gradient-removal principle; compare
\cite[Section~5, Lemma~6]{ImbertSilvestre2013}.  We give the local
proof for functions allowed to take the value $-\infty$, rather than
apply that lemma directly.

\begin{lemma}[Removal of the exceptional gradient]
\label{lem:exceptional-gradient}
Let $W\colon\Omega\to[-\infty,\infty)$ be upper semicontinuous and
locally bounded above on a connected planar domain.  Suppose every
smooth function $\psi$ touching $W$ from above at a finite point
$x_0$, with $\nabla\psi(x_0)\ne0$, satisfies
$\Delta\psi(x_0)\ge0$.  Then $W$ is subharmonic, with the
identically $-\infty$ function allowed.
\end{lemma}

\begin{proof}
Suppose first that an upper test at $0$ has zero gradient and
negative Laplacian; translate the point to $0$ if necessary.  A
quadratic Taylor polynomial with its Hessian increased by a
sufficiently small positive multiple of the identity gives a
quadratic upper test $h$ with
\[
 \begin{gathered}
 W(0)=h(0),\qquad \nabla h(0)=0,\qquad
 \operatorname{tr}h''<0,\\
 W(x)<h(x)\quad(0<|x|\le r)
 \end{gathered}
\]
for some small $r>0$.  Choose a unit eigenvector $e$ of $h''$ with
eigenvalue $\alpha<0$.  For small $t>0$, maximize
$W(x)-h(x)+t e\cdot x$ on $\overline{\D_r}$.  Upper
semicontinuity and the strict boundary gap ensure that a maximizer
$x_t$ is interior.  Its value is at least the value $0$ at the
origin, and $W-h\le0$, so $e\cdot x_t\ge0$.  At $x_t$ an upper
test is $h(x)-t e\cdot x$, shifted by its contact constant.  Its
gradient has component
\[
 e\cdot\nabla(h-t e\cdot x)(x_t)
 =\alpha\, e\cdot x_t-t<0,
\]
while its Laplacian is still negative.  This contradicts the
hypothesis.  All smooth upper tests at finite points therefore
have nonnegative Laplacian.

We obtain harmonic comparison directly.  Let
$\overline{\D_r(a)}\subset\Omega$, let $h$ be harmonic in
$\D_r(a)$ and continuous on its closure, and suppose $W\le h$ on
the boundary.  If $W-h$ were positive somewhere inside, then for
sufficiently small $\eps>0$ the function
\[
 h_\eps(x)=h(x)+\eps(r^2-|x-a|^2)
\]
would still lie below $W$ at some interior point.  A positive
maximum of $W-h_\eps$ is attained at an interior finite point.
Adding this maximum to $h_\eps$ supplies an upper test with
Laplacian $-4\eps<0$, again a contradiction.  Thus $W$ satisfies
harmonic comparison.  For upper semicontinuous functions this is
the usual characterization of subharmonicity, including the
identically $-\infty$ case; see \cite[Chapter~2]{Ransford1995}.
In particular, local integrability
and the submean property of a nontrivial subharmonic function are
consequences here, not assumptions used in the argument.
\end{proof}

\begin{proof}[Proof of Theorem~\ref{thm:subharmonic-profile}]
Fix such a profile $V$ and pass to a further subsequence so
that $S_j/D_j\to d$ and $B_j/D_j\to b$ as in
Lemma~\ref{lem:carleman}.  The profile bounds remain valid on this
subsequence.  Suppose a smooth $\chi$ touches $V$ from above at
a finite point $x_0$, and suppose
\[
 \Delta\chi(x_0)<0,
 \qquad b+d\nabla\chi(x_0)\ne0.
\]
By adding a small positive multiple of $|x-x_0|^4$, we make the
touching strict on a sufficiently small closed disk without
altering the gradient or Hessian at $x_0$.  Take this disk inside
the neighborhood in Lemma~\ref{lem:carleman}.

Choose fixed radii $0<r_0<r_1<r_2<r_3$ inside that neighborhood.
Let $\zeta$ be a smooth cutoff equal to one on $\D_{r_1}(x_0)$
and supported in $\D_{r_2}(x_0)$.  Its derivatives are supported
in the annulus between these two disks.  Set
\[
 G=\D_{r_0/2}(x_0),\qquad
 F=\{x:r_0\le |x-x_0|\le r_3\}.
\]
We may choose the radii so that the strict touching holds on the
whole closed outer disk.  Then $F$ is a compact annulus,
$\sup_F(V-\chi)<0$, and $\sup_G(V-\chi)=0$.

For completeness, the cutoff error can be controlled with a
prefactor independent of the ratio $D_j/S_j$.  Choose a smooth
$\theta$ supported in the interior of $F$ and equal to one near
the derivatives of $\zeta$.  Testing the equation for $U_j$
against $\theta^2e^{2T_j}U_j$ and integrating by parts bounds
\[
 \begin{aligned}
 \int\theta^2e^{2T_j}|\nabla U_j|^2
 &\le C\int_F e^{2T_j}
       \bigl(|\nabla\theta|^2+
       \theta^2|\nabla T_j|^2+K_j^2|p_j|\theta^2\bigr)|U_j|^2\\
 &\le CD_j^2\int_F e^{2T_j}|U_j|^2.
 \end{aligned}
\]
The first inequality absorbs half of the gradient integral by
Cauchy--Schwarz; the second uses $K_j=o(D_j)$ and $D_j\to\infty$.
Since
\[
 (\Delta+K_j^2p_j)(\zeta U_j)
   =(\Delta\zeta)U_j+2\nabla\zeta\cdot\nabla U_j,
\]
Lemma~\ref{lem:carleman} gives
\[
 c\sqrt{S_j}D_j\|e^{T_j}U_j\|_{L^2(G)}
 \le CD_j\|e^{T_j}U_j\|_{L^2(F)}.
\]
Cancel $D_j$ before comparing exponential rates.  Squaring and
multiplying by the common normalization constant $c_j$ yields
\[
 S_j\int_G e^{-2S_j\chi}\,d\mu_j
       \le C\int_F e^{-2S_j\chi}\,d\mu_j.
\]
The weighted bounds of Lemma~\ref{lem:profile-extraction}, together
with $\log S_j/S_j\to0$, now give the contradiction
\[
 0=\sup_G(V-\chi)\le\sup_F(V-\chi)<0.
\]
In particular, no bound on $\log D_j/S_j$ was required.

We have established the upper-test inequality whenever
$b+d\nabla\chi(x_0)\ne0$.  If $d=0$, then $|b|=1$ and there is
no exceptional gradient; harmonic comparison follows as in
Lemma~\ref{lem:exceptional-gradient}.  If $d>0$, put
$W(x)=V(x)+(b/d)\cdot x$.  An upper test $\psi$ for $W$ corresponds
to the upper test $\chi=\psi-(b/d)\cdot x$ for $V$, and
$b+d\nabla\chi=d\nabla\psi$.  Thus $W$ satisfies the hypothesis
of Lemma~\ref{lem:exceptional-gradient}.  It is subharmonic, and
subtracting the affine function shows that $V$ is subharmonic too.
\end{proof}

The next estimate is a fixed-disk form of quantitative unique
continuation.  For the classical doubling estimate of Donnelly and
Fefferman and the corresponding propagation bound, see
\cite[Theorem~4.1 and Lemma~4.14]{NazarovPolterovichSodin2005}.
We give a proof using logarithmic profiles, which also supplies the
local nonvanishing needed below.

\begin{proposition}[Initial growth on a fixed disk]
\label{prop:initial-growth}
Let $(M,g)$ be a fixed smooth closed connected surface, and let
$G\subset M$ be a fixed nonempty open coordinate disk.  There is
a finite constant $C=C(M,g,G)$ such that every nonzero real
eigenfunction $-\Delta_g u=k^2u$ with $k>0$ satisfies
$\|u\|_{L^2(G)}>0$ and
\begin{equation}\label{eq:initial-growth}
 \log\frac{\|u\|_{L^2(M)}}{\|u\|_{L^2(G)}}\le Ck.
\end{equation}
Both norms use Riemannian volume.  The constant is uniform over
all eigenfunctions in every eigenspace.
\end{proposition}

\begin{proof}
The positive spectral gap gives $k\ge k_1>0$.  Interpret the log
ratio in Equation~\eqref{eq:initial-growth} as $+\infty$ when the
local norm is zero.  If the assertion failed, we could choose
eigenfunctions $u_j$, normalized by $\|u_j\|_{L^2(M)}=1$, such
that their log ratios $R_j$ obey $R_j\ge j^2 k_j$.  This also
covers a zero local norm by repeating that eigenfunction.  Set
$S_j=j k_j$.  Then $S_j\to\infty$ and $S_j/k_j\to\infty$,
whether or not $k_j$ tends to infinity.

Apply Lemma~\ref{lem:profile-extraction} on $M$ to the probability
measures $|u_j|^2\,d\operatorname{vol}_g$.  Compactness of $M$ and
the compact upper bound with $F=M$ give
$0\le\sup_M V\le0$.  Upper semicontinuity ensures that this
maximum is attained.

In a relatively compact isothermal chart,
$g=p(x)(dx_1^2+dx_2^2)$ with $p>0$ smooth, and the coordinate
functions satisfy $\Delta U_j+k_j^2pU_j=0$.  The positive
volume factor $p$ is bounded above and below on this chart.
Multiplying $|U_j|^2\,dx$ by a fixed positive constant if
necessary makes its mass at most one; the resulting logarithmic
profile is the restriction of $V$.  Theorem~\ref{thm:subharmonic-profile}
applies with $B_j=0$, $K_j=k_j$, and $D_j=S_j$, since
\[
 \frac{k_j}{S_j}=\frac1j\longrightarrow0,
 \qquad
 \frac{k_j^2\|\nabla p\|_\infty}{S_j^2}
       =\frac{\|\nabla p\|_\infty}{j^2}\longrightarrow0.
\]
Thus $V$ is subharmonic in every such chart.  The strong maximum
principle makes the set where $V=0$ open; it is also closed by
upper semicontinuity and nonempty by the preceding maximum
argument.  Connectedness gives $V\equiv0$ on $M$.

On the other hand, the normalized mass of the given disk satisfies
\[
 \frac{1}{2S_j}\log\int_G |u_j|^2\,d\operatorname{vol}_g
   =-\frac{R_j}{S_j}\le-j,
\]
with the same conclusion when the integral is zero.  This
contradicts the open lower bound in
Equation~\eqref{eq:profile-bounds}, whose right-hand side is
$\sup_GV=0$.  The argument proves both positivity of the local
norm and Equation~\eqref{eq:initial-growth}.  It uses only local
isothermal charts and therefore imposes no orientability
assumption.
\end{proof}
